\chapter{Conclusion and Future Scope}

\section{Summary of the Study}

This project presented \textbf{Smart Interview AI}, a full-stack, AI-powered mock interview platform designed to democratise access to high-quality interview preparation. The platform was built using the MERN stack (MongoDB, Express.js, React, Node.js) augmented by a Python FastAPI AI server and the Google Gemini 2.5 Flash large language model.

The system successfully delivers on its core objectives:
\begin{itemize}
    \item \textbf{Personalised question generation:} The Gemini API generates role-specific, difficulty-calibrated interview questions from the candidate's resume and target role, covering five interview types.
    \item \textbf{Real-time speech transcription:} The Web Speech API transcribes spoken answers in real time, enabling a natural interview experience without manual typing.
    \item \textbf{Coding interview support:} A Monaco-based code editor with Piston API execution supports live coding challenges across 13 programming languages, with robust argument normalisation for Python and JavaScript.
    \item \textbf{AI-generated feedback:} After each session, Gemini generates a comprehensive feedback report covering content metrics, strengths, improvements, and personalised recommendations.
    \item \textbf{Resume analysis:} The platform parses uploaded resumes, extracts structured information, and computes an ATS compatibility score.
    \item \textbf{Production deployment:} The frontend is deployed on Vercel and the backend on Render, with MongoDB Atlas as the cloud database.
\end{itemize}

\section{Main Findings}

\begin{enumerate}
    \item \textbf{LLM-based question generation is effective:} Gemini 2.5 Flash consistently generates contextually relevant, role-specific questions. The quality is comparable to questions asked in real technical interviews, particularly for software engineering roles.

    \item \textbf{Background analysis races with session termination:} A critical architectural issue was identified where background AI analysis of responses could be incomplete when the user ends the interview. The fix---running synchronous analysis during feedback generation if metrics are zero---resolves this reliably.

    \item \textbf{Argument normalisation is essential for code execution:} Users write functions with arbitrary names and parameter orders. A robust argument normalisation pipeline using Python introspection and a fallback calling chain is necessary to achieve high test case pass rates.

    \item \textbf{Free-tier deployment has significant constraints:} Render's free tier (512MB RAM) cannot run the Python AI server with MediaPipe and Librosa. The TypeScript compilation must occur during the build phase, not at runtime, to avoid OOM crashes.

    \item \textbf{DNS configuration is critical for MongoDB Atlas:} The \texttt{querySrv} DNS lookup for \texttt{mongodb+srv} URIs fails on some ISP DNS servers. Overriding Node.js DNS to use Google's servers (8.8.8.8) resolves this reliably.
\end{enumerate}

\section{Contributions to Knowledge}

\begin{itemize}
    \item Demonstrated the feasibility of building a production-grade, AI-powered interview simulation platform using entirely free-tier cloud services.
    \item Developed a robust Python test harness that normalises multi-line, multi-type function arguments for automated code evaluation.
    \item Identified and resolved a race condition between background AI analysis and interview session termination in a fire-and-return API pattern.
    \item Documented a comprehensive set of deployment issues and fixes for MERN stack applications on Render and Vercel.
\end{itemize}

\section{Practical Implications}

Smart Interview AI has direct practical value for:
\begin{itemize}
    \item \textbf{Job candidates:} Provides affordable, on-demand interview practice with personalised feedback, reducing dependence on expensive coaching services.
    \item \textbf{Educational institutions:} Can be integrated into placement preparation programmes to give students structured interview practice.
    \item \textbf{Recruiters:} The platform's question generation and evaluation capabilities could be adapted for automated pre-screening of candidates.
    \item \textbf{Developers:} The codebase serves as a reference implementation for integrating LLMs, real-time communication, and computer vision in a production web application.
\end{itemize}

\section{Future Work}

The following enhancements are planned for future development:

\begin{enumerate}
    \item \textbf{Python AI Server Deployment:} Deploy the MediaPipe-based video analysis and Librosa audio analysis server on Railway or Fly.io after stripping unused heavy dependencies (\texttt{torch}, \texttt{transformers}).

    \item \textbf{Interview Recording:} Implement video blob upload to Cloudinary and playback on the feedback page, enabling candidates to review their performance.

    \item \textbf{Real PDF Export:} Replace \texttt{window.print()} with \texttt{jspdf} + \texttt{html2canvas} for generating downloadable PDF feedback reports.

    \item \textbf{WebRTC with TURN Server:} Add a TURN server (Twilio or coturn) to enable reliable peer-to-peer video connections behind NAT for future multi-participant interview modes.

    \item \textbf{AssemblyAI Integration:} Replace the rate-limited Google Speech Recognition fallback with AssemblyAI for production-grade, server-side speech transcription.

    \item \textbf{OTP Two-Factor Authentication:} Implement TOTP-based 2FA using \texttt{speakeasy} for enhanced account security.

    \item \textbf{Automated Scheduling Reminders:} Add a \texttt{node-cron} job to automatically send email reminders for scheduled interviews.

    \item \textbf{Cross-Interview Improvement Tracking:} Implement comparison of current interview scores against previous sessions to show improvement trends on the feedback page.

    \item \textbf{Mobile Application:} Develop a React Native mobile application for iOS and Android.

    \item \textbf{Multi-Language Support:} Add support for interviews in Hindi and other regional languages using multilingual LLM prompting.
\end{enumerate}

\section{Final Thoughts}

Smart Interview AI demonstrates that modern AI capabilities---large language models, computer vision, and speech recognition---can be combined in a unified, accessible web platform to meaningfully improve the interview preparation experience. The project successfully navigated significant technical challenges including LLM integration, real-time multimodal analysis, robust code execution, and cloud deployment constraints.

The platform is production-ready for its core features (interview simulation, coding challenges, and AI feedback) and provides a solid foundation for the planned enhancements. As LLM capabilities continue to improve and cloud infrastructure costs decrease, AI-powered interview preparation platforms like Smart Interview AI will become an increasingly important tool for candidates worldwide.
